\section{Sorting and Searching from CSES}

\Topic{Technique map}
The most important decision is which \emph{shape} the search space has:
\begin{itemize}
	\item \textbf{Sort + scan:} For multiset properties. (Distinct Numbers, Restaurant Customers)
	\item \textbf{Two pointers:} For monotonic boundaries. (Sum of Two Values, Apartments)
	\item \textbf{Binary search on answer:} For monotone predicates. (Factory Machines)
	\item \textbf{Dynamic sorted set:} For predecessor/successor queries. (Concert Tickets, Traffic Lights)
	\item \textbf{Prefix sums + ordered structure:} For subarray constraints. (Subarray Sums)
\end{itemize}

\Topic{Sum of Two / Three / Four Values}

\Problem Find distinct positions whose values sum to $x$.

Keep indices before sorting. Two-sum is sorted two-pointers; three-sum fixes one; four-sum uses a map.
\begin{lstlisting}
// two values
sort (value,index)
i=0, j=n-1
while i<j:
    s=a[i]+a[j]
    if s==x: answer
    if s<x: i++ else j--

// four values
for i=0..n-1:
    for j=i+1..n-1:
        if x-a[i]-a[j] in oldPairs: answer
    for j=0..i-1:
        oldPairs[a[i]+a[j]]=(i,j)
\end{lstlisting}

\Solution Two-sum uses a two-pointer sweep on sorted data. Four-sum stores pair sums in a hash map, adding to the map only indices $<i$ to ensure disjointness.

\Analysis Staggering insertion in four-sum avoids overlapping indices. Fixing an element for three-sum reduces it to two-sum.

\Topic{Playlist}

\Problem Find the longest contiguous segment of distinct songs.

Store the last position of each song and move the left boundary past duplicates.
\begin{lstlisting}
l=0
for r=0..n-1:
    if seen a[r]: l=max(l,last[a[r]]+1)
    last[a[r]]=r
    ans=max(ans,r-l+1)
\end{lstlisting}

\Solution A sliding window $[l,r]$ avoids duplicates by jumping $l$ past the last seen index of $a[r]$. $O(n)$ time using a hash map.

\Analysis The \texttt{max} ensures $l$ never moves backward when an old duplicate is found outside the current window.

\Topic{Concert Tickets}

\Problem Sell the most expensive affordable ticket to each customer in order.

Sell the largest ticket not exceeding each customer's price.
\begin{lstlisting}
multiset tickets
for x in customers:
    it = tickets.upper_bound(x)
    if it == begin: print -1
    else:
        --it
        print *it
        erase it
\end{lstlisting}

\Solution Use a multiset to query the predecessor of the customer's budget. Greedily taking the largest affordable ticket saves cheaper tickets. $O((n+m) \log n)$ time.

\Analysis Predecessor queries with deletion are perfect for balanced BSTs (multisets).

\Topic{Restaurant Customers}

\Problem Find the maximum overlapping customer intervals.

Sweep arrival/leave events to find maximum overlap.
\begin{lstlisting}
events += (arrival,+1), (leave,-1)
sort by time; process leave before arrival on tie
cur=ans=0
for event: cur += delta, ans=max(ans,cur)
\end{lstlisting}

\Solution Event sweep converts interval overlap into a prefix sum. Sort by time and track the running total. $O(n \log n)$ time.

\Analysis Processing departures before arrivals on ties handles CSES's non-overlapping endpoint semantics correctly.

\Topic{Stick Lengths}

\Problem Change $n$ stick lengths to be equal, minimizing $\sum |a_i - x|$.

The best target length is the median.
\begin{lstlisting}
sort a
m=a[n/2]
ans=sum(abs(x-m))
\end{lstlisting}

\Solution The median minimizes the sum of absolute differences. $O(n \log n)$ time.

\Analysis The cost function is convex. Moving the target shifts the penalty equally between left and right elements, zeroing out at the median.

\Topic{Towers}

\Problem Stack cubes. New cubes must be strictly smaller than the tower top. Minimize towers.

Patience sorting: replace the smallest tower top strictly greater than $x$.
\begin{lstlisting}
multiset tops
for x in cubes:
    it=tops.upper_bound(x)
    if it exists: erase it
    tops.insert(x)
answer=tops.size()
\end{lstlisting}

\Solution Greedily placing a cube on the tightest valid tower top preserves larger tops for future cubes. $O(n \log n)$ time.

\Analysis This corresponds to patience sorting. By Dilworth's theorem, towers form the longest non-decreasing subsequence.

\Topic{Traffic Lights}

\Problem Add lights to a street and report the longest gap after each addition.

Maintain light positions and segment lengths.
\begin{lstlisting}
pos={0,x}, len={x}
for p in insertions:
    r=*pos.upper_bound(p), l=*prev(...)
    erase len r-l
    insert len p-l and r-p
    pos.insert(p)
    print max(len)
\end{lstlisting}

\Solution A set tracks positions and a multiset tracks gap lengths. Inserting a light splits one gap into two. $O(n \log n)$ time.

\Analysis Tracking both endpoints and gap lengths explicitly allows $O(\log n)$ updates and maximum queries.

\Topic{Factory Machines}

\Problem Find the minimum time for $n$ machines to make $k$ products.

Binary search the minimum time.
\begin{lstlisting}
ok(T):
    made=0
    for machine in a:
        made += T/machine
        if made >= need: return true
    return false

lo=0, hi=1e18
while lo<hi:
    mid=(lo+hi)/2
    if ok(mid): hi=mid
    else: lo=mid+1
\end{lstlisting}

\Solution Production capacity is monotonic with time. Binary search the answer up to $10^{18}$. $O(n \log (\text{max time}))$ time.

\Analysis The early exit inside \texttt{ok()} prevents integer overflow when testing large $T$.

\Topic{Array Division}

\Problem Divide array into $k$ contiguous subarrays to minimize maximum sum.

Minimize maximum segment sum using binary search + greedy feasibility.
\begin{lstlisting}
ok(limit):
    groups=1, cur=0
    for x in a:
        if x>limit: return false
        if cur+x>limit:
            groups++, cur=0
        cur += x
    return groups <= k
\end{lstlisting}

\Solution Binary search the maximum allowed sum. The greedy check forms the fewest groups by packing subarrays as tightly as possible. $O(n \log (\sum a_i))$ time.

\Analysis The monotonic search space ranges from $\max(a_i)$ to $\sum a_i$.

\Topic{Subarray Sums and Maximum Subarray Sum II}

\Problem Count subarrays with exact sum (Sums I/II), or find max subarray sum within length bounds.

Positive values allow sliding window. Arbitrary values require prefix sums.
\begin{lstlisting}
// Subarray Sums II
cnt[0]=1
pref=ans=0
for x in a:
    pref += x
    ans += cnt[pref-target]
    cnt[pref]++

// max sum with length [A,B]
insert pref[r-A]
erase pref[r-B-1]
ans=max(ans, pref[r]-min(active))
\end{lstlisting}

\Solution For exact sum, hash map counts past prefix sums matching the needed difference. For bounded max sum, a multiset queries the minimum valid left prefix.

\Analysis Prefix sums convert subarray constraints into differences between two points, reducing to point lookups or sliding window extremum queries.
